\documentclass{article}
\input{../../lib.sty}

\title{\texttt{impl Sample for ZExpFamily<2>}}
\author{Michael Shoemate}
\date{}

\begin{document}

\maketitle

\contrib
Proves soundness of the implementation of \rustdoc{measurements/noise/trait}{Sample} for \texttt{ZExpFamily<2>} in \asOfCommit{mod.rs}{f5bb719}.

\section{Hoare Triple}
\subsection*{Precondition}
\texttt{self} represents a valid probability distribution.

\subsection*{Pseudocode}
\lstinputlisting[language=Python,firstline=2,escapechar=|]{./pseudocode/Sample_for_ZExpFamily2.py}

\subsection*{Postcondition}
\begin{theorem}
    \label{postcondition}
    Either returns \texttt{Err(e)} independently of the input \texttt{shift},
    or returns a sample from the additive mechanism defined by \texttt{self} shifted by
    \texttt{shift}. When \texttt{self.divisor} is \texttt{None}, this is \texttt{shift + Z}.
    When \texttt{self.divisor = m}, this is \texttt{(shift + Z) mod m}.
\end{theorem}

\begin{proof}
    By the precondition, since \texttt{self} represents a valid probability distribution, 
    then by the definition of \rustdoc{measurements/noise/struct}{ZExpFamily},
    \texttt{self.scale} is non-negative.

    Since the preconditions for \rustdoc{traits/samplers/cks20/sample\_discrete\_gaussian} are met (non-negative scale),
    the sample drawn on line \ref{line:add} has the base discrete-Gaussian law encoded by \texttt{self.scale}.

    If \texttt{self.divisor} is absent, line \ref{line:add} returns \texttt{shift + Z}, establishing
    the non-modular case.
    If \texttt{self.divisor} is present, line \ref{line:mod} deterministically reduces the result modulo
    the divisor, which is exactly the wrapped additive mechanism encoded by \texttt{self}.
    Therefore the implementation satisfies the postcondition in either branch.
\end{proof}

\end{document}
